\chapter{The Mobile Application}
\label{ch:mobile}

The mobile client brings \voname{}'s messaging to phones. It is a Flutter
application focused on the chat domain --- channels, direct messages, threads,
read receipts, typing indicators, and, importantly, push notifications that work
whether the app is in the foreground, the background, or fully terminated. This
chapter describes its clean-architecture layering, its Provider-based state
management, its dual-transport networking, its real-time and notification stacks,
and the engineering details that make it feel responsive and reliable.

\section{Overview and Scope}

The application is built with Flutter and Dart and uses a small, deliberate set of
packages: \textbf{Provider} for state management, \code{http} for REST,
\textbf{\code{stomp\_dart\_client}} for real-time messaging,
\textbf{\code{awesome\_notifications}} for local/push notifications,
\code{shared\_preferences} for persistence, and \code{uuid} for client-generated
message ids. Its scope is intentionally the \emph{chat} experience: the mobile app
does not implement the 2D office, workspace administration, or in-app calling.
This focus lets it implement the messaging domain to a high standard --- optimistic
delivery, reconnection catch-up, and cross-state notifications --- rather than
implementing everything shallowly.

\section{Architecture}

The code follows a \textbf{feature-first, clean-architecture-inspired} layering
in which dependencies flow strictly downward (Figure~\ref{fig:mob-layers}). The
\emph{repository} is the linchpin: it is the single source of truth that merges
the REST data source, the WebSocket data source, and an in-memory cache, and
re-broadcasts consolidated state through its own streams that the providers
consume.

\begin{figure}[htbp]
\centering
\begin{tikzpicture}[font=\footnotesize,
  layer/.style={draw=voPrimary,fill=voPrimary!8,rounded corners=3pt,minimum height=0.85cm,
                text width=10.5cm,align=center},
  a/.style={-{Stealth[]},thick,voGray}]
\node[layer] (ui) at (0,4) {\textbf{Presentation} --- screens \& widgets (Consumer / context.watch)};
\node[layer] (st) at (0,3) {\textbf{State} --- Providers (ChangeNotifier): Connection, Channel, Message, Thread};
\node[layer] (repo) at (0,2) {\textbf{Repository} --- single source of truth: merges REST + WS + cache, owns notification logic};
\node[layer,text width=5.1cm] (rest) at (-2.7,0.8) {\textbf{REST data source}\\ChatApiClient (http)};
\node[layer,text width=5.1cm] (ws) at (2.7,0.8) {\textbf{WS data source}\\StompService};
\node[layer] (mod) at (0,-0.2) {\textbf{Models} --- DTOs with fromJson / toJson};
\draw[a] (ui)--(st); \draw[a] (st)--(repo);
\draw[a] (repo)--(rest); \draw[a] (repo)--(ws);
\draw[a] (rest)--(mod); \draw[a] (ws)--(mod);
\end{tikzpicture}
\caption{Mobile clean-architecture layering; data flows WS/REST $\to$ repository
$\to$ providers $\to$ widgets.}
\label{fig:mob-layers}
\end{figure}

Dependency injection is manual, through Provider's \code{MultiProvider}: a single
root \code{ConnectionProvider} lazily constructs the repository (and its REST and
WebSocket data sources) on a successful connection, and hands it down to
feature-scoped providers created as screens are pushed.

\section{State Management with Provider}

Five \code{ChangeNotifier} providers structure the application's state:

\begin{description}
  \item[\code{ConnectionProvider}] (root) owns the server configuration and
  connection status and the singleton repository. Its \code{connect()} normalises
  the URLs, runs a health check, opens the socket, loads the workspace's channels
  and DMs, and \emph{auto-subscribes to all of them} so that notifications arrive
  app-wide even for conversations the user is not currently viewing.
  \item[\code{ChannelProvider}] holds the channel and DM lists and per-channel
  unread counts, and exposes load/join/create/leave operations.
  \item[\code{MessageProvider}] manages one open conversation: a
  \code{createdAt}-sorted message list, reply/edit state, and a debounced typing
  indicator; it subscribes to the repository's message and typing streams filtered
  to its channel and performs catch-up on reconnect.
  \item[\code{ThreadProvider}] and \code{ThreadMessageProvider} mirror the channel
  and message providers for threads.
\end{description}

A recurring correctness theme is safe disposal: providers guard against calling
\code{notifyListeners} after disposal and cache the repository reference so that
teardown never touches a \code{BuildContext} that may already be unmounted.

\section{Networking: A Dual-Transport Repository}

The repository speaks to the chat service over two transports and degrades
gracefully between them.

\paragraph{REST.} The \code{ChatApiClient} wraps the \code{http} package with the
identity headers (\code{X-User-Id}, \code{X-User-Role}), a ten-second timeout, and
typed error handling. It covers the full chat surface: channels and DMs, message
history (with \code{before}/\code{after} cursors for pagination and catch-up),
send/edit/delete, threads, and read receipts.

\paragraph{Real-time.} The \code{StompService} connects to
\route{/api/chat/connect?ticket=<t>} (the same ticket pattern as the web client),
auto-reconnects, and exposes the incoming message and typing events as Dart
streams. It subscribes to \route{/topic/channel/\{id\}} and its typing sibling,
and to the thread equivalents, and publishes sends to \route{/app/chat/send}.

\paragraph{Optimistic delivery.} Sending is optimistic
(Figure~\ref{fig:mob-send}): the repository immediately inserts a temporary
message keyed by a generated \code{clientMessageId}, attempts delivery over the
WebSocket, and \emph{falls back to a REST \code{POST}} if the socket send throws.
When the server's authoritative echo arrives --- matched by that same
\code{clientMessageId} --- it replaces the placeholder; on total failure the
optimistic message is rolled back. On reconnection the repository fetches any
messages that arrived while it was offline (everything \code{after} the last known
id), so no message is missed.

\begin{figure}[htbp]
\centering
\begin{tikzpicture}[font=\footnotesize,
  b/.style={draw=voPrimary,fill=voPrimary!8,rounded corners=2pt,minimum height=0.7cm,
            text width=2.3cm,align=center},
  a/.style={-{Stealth[]},thick,voDark},
  f/.style={-{Stealth[]},thick,voAccent!60!black,densely dashed},
  note/.style={font=\scriptsize\itshape,voGray}]
\node[b] (ui) at (0,0) {UI};
\node[b] (rp) at (3.5,0) {Repository};
\node[b] (ws) at (7.2,0.8) {WebSocket};
\node[b] (rs) at (7.2,-0.8) {REST};
\draw[a] (ui) -- node[above,note]{send} (rp);
\node[note,align=center] at (3.5,-0.7) {insert tmp\_<id>};
\draw[a] (rp) -- node[above,note]{try first} (ws);
\draw[f] (rp) -- node[below,note]{fallback} (rs);
\draw[a] (ws) to[out=100,in=60] node[above,note]{echo (clientMessageId) reconciles} (rp.60);
\end{tikzpicture}
\caption{Optimistic send with WebSocket-first, REST-fallback, and echo
reconciliation.}
\label{fig:mob-send}
\end{figure}

\section{Identity on Mobile}

The mobile app takes a pragmatic approach to identity suited to a development and
demonstration client: rather than a full login screen, a \emph{connection screen}
lets the user enter the server URLs, a numeric user id, and a role, which are
persisted with \code{shared\_preferences} and sent as the \code{X-User-Id} and
\code{X-User-Role} headers --- exactly the headers the gateway would otherwise
inject after validating a JWT. The WebSocket still uses the proper single-use
ticket handshake. This mirrors, on the client, the platform's server-side trust
model (Chapter~\ref{ch:architecture}); wiring a real JWT login is a
straightforward extension, since the backend already issues and validates the
token.

\section{Notifications Across App States}

The notification system is the mobile client's most sophisticated piece. Built on
\code{awesome\_notifications}, it defines four styled channels (channel messages,
direct messages, thread replies, and system notices) and, critically, registers
its lifecycle handlers as static, entry-point-annotated functions so they execute
even when the app is \emph{backgrounded or terminated}.

The decision of \emph{whether} to raise a notification lives in the repository: a
notification fires only when a message is not from the current user \emph{and} the
user is not currently viewing that channel or thread (tracked by ``active''
sets updated as screens are entered and left). Each notification embeds a payload
(channel id, thread id, DM flag, message id) so that tapping it deep-links to the
right conversation. A small navigation bridge makes this reliable across all three
app states: if the app is running, it navigates immediately; if it was terminated,
the target route is stored and replayed once the app's navigator is ready on the
first frame.

\section{Data Models and Bootstrapping}

The model layer is a set of null-safe DTOs with \code{fromJson}/\code{toJson} ---
\code{ChannelResponse}, \code{MessageResponse} (with \code{deleted}, \code{threadId},
\code{replyToId}, \code{mentions}, and \code{clientMessageId}),
\code{ThreadResponse}, a generic \code{PaginatedResponse}, and the typing and
error envelopes --- mirroring the chat service's contract. Application startup
initialises the notification service (its channels, permissions, and background
handlers) \emph{before} the app runs, then mounts the root provider and a
\code{MaterialApp} that holds a global navigator key (so the notification bridge
can navigate from outside the widget tree) and routes to the connection screen.

\section{Engineering Highlights}

For the evaluation in Chapter~\ref{ch:testing}, the mobile client demonstrates
several production-grade patterns worth naming: optimistic UI with
\code{clientMessageId} reconciliation; dual-transport delivery with graceful
degradation; reconnection resilience with message catch-up; disposal safety; and
correct notification behaviour across foreground, background, and terminated
states. Its main limitation is one of \emph{breadth} rather than depth: it covers
only the chat domain, leaving workspace administration, the 2D office, and in-app
calling to the web and game clients.
